\section{État de l'Art}

\subsection{Méthodologie de la revue de littérature}

Cette revue de l'état de l'art a été conduite en croisant deux types de sources, conformément aux exigences
méthodologiques d'un mémoire de licence en intelligence artificielle et data : des sources bibliographiques
scientifiques (articles évalués par les pairs, prépublications arXiv, actes de conférences ACM/IEEE,
chapitres Springer) et des sources webographiques (blogs techniques spécialisés, dépôts GitHub de
référence, livres blancs d'éditeurs). Les bases consultées incluent arXiv, Springer Link, ScienceDirect et
ResearchGate. Les mots-clés de recherche combinaient les trois axes techniques du projet : retrieval-augmented
generation, analyse de sentiment par modèles de langage, et web scraping appliqué à l'intelligence
concurrentielle.

\subsection{Études bibliographiques}

\subsubsection{Retrieval-Augmented Generation (RAG)}

Le paradigme RAG a été introduit par Lewis et al. (2020) \cite{lewis2020} dans un article fondateur présenté
à NeurIPS, qui propose de coupler un modèle génératif à un mécanisme de récupération d'information externe
afin de pallier les limites des modèles de langage purement paramétriques : hallucinations, obsolescence des
connaissances et absence d'expertise de domaine spécifique. Cette approche a depuis fait l'objet de
nombreuses synthèses. Wu et al. (2024) \cite{wu2024} proposent une revue systématique centrée sur les
mécanismes de récupération et de fusion (query-based, logits-based, latente, paramétrique), et établissent
une taxonomie comparant ces approches selon leur accessibilité, leur efficacité et leurs cas d'usage. Yu et
al. (2024) \cite{yu2024} se concentrent quant à eux sur l'évaluation des systèmes RAG et proposent un cadre
unifié (Auepora) structurant les métriques autour de trois axes : la pertinence de la récupération,
l'exactitude de la génération et la fidélité aux sources récupérées.

Plusieurs travaux récents traitent des limites intrinsèques du RAG de base et proposent des correctifs
pertinents pour ce projet. Yan et al. (2024) \cite{yan2024} introduisent le Corrective RAG, qui ajoute une
étape de vérification de la pertinence des documents récupérés avant leur transmission au modèle générateur
--- une logique proche du seuil de confiance mis en œuvre dans ce projet. Xu et al. (2024) \cite{xu2024},
lors de SIGIR 2024, montrent l'intérêt de coupler RAG et graphes de connaissances pour des cas d'usage de
service client, un contexte proche de l'assistance commerciale visée ici. Xue et al. (2024) \cite{xue2024}
documentent par ailleurs des vulnérabilités de sécurité propres aux systèmes RAG (BadRAG), rappelant que la
fiabilité de ces architectures reste un sujet de recherche actif plutôt qu'un problème résolu.

\subsubsection{Analyse de sentiment par grands modèles de langage}

L'analyse de sentiment a longtemps reposé sur des modèles supervisés entraînés spécifiquement sur la tâche
(SVM, puis modèles Transformer de type BERT ou FinBERT). L'émergence des grands modèles de langage
généralistes a introduit un nouveau paradigme, le zero-shot, où le modèle classe le sentiment sans exemple
d'entraînement spécifique. Simmering et Huoviala (2023) \cite{simmering2023} comparent GPT-4 et GPT-3.5 en
zero-shot, few-shot et fine-tuné sur une tâche d'analyse de sentiment par aspect (ABSA) : le modèle fine-tuné
atteint un score F1 état de l'art de 83,8, mais au prix d'un coût d'inférence largement supérieur ; les
auteurs concluent que des prompts détaillés améliorent significativement les performances en zero-shot et
few-shot, rendant cette approche viable en l'absence de données annotées.

Cette conclusion est corroborée par une étude publiée dans JMIR Formative Research (2025) évaluant GPT-4 sur
des données de santé publique : la précision en configuration zero-shot atteint déjà 92\,\% pour le sentiment
positif, avec une amélioration marginale seulement en few-shot, ce qui suggère que les grands modèles
récents disposent d'une capacité de classification de sentiment déjà robuste sans exemples. Vamvourellis et
Mehta (2025) \cite{vamvourellis2025}, dans une étude menée chez BlackRock sur le sentiment financier,
comparent plusieurs LLM propriétaires en zero-shot face à des modèles FinBERT spécialisés fine-tunés, et
notent que les stratégies de raisonnement explicite (chain-of-thought) n'améliorent pas systématiquement
l'alignement avec l'annotation humaine, en particulier sur les cas peu ambigus. Une synthèse plus large
portant sur 13 tâches et 26 jeux de données (citée dans une revue de 2026 sur l'analyse de sentiment
appliquée aux critiques de films) confirme que les LLM en zero-shot restent compétitifs sur des tâches de
classification simple, mais restent en retrait des modèles spécialisés fine-tunés sur les tâches nécessitant
une compréhension fine de phénomènes de sentiment complexes. Beno (2025) \cite{beno2025} montre enfin qu'une
approche hybride combinant un modèle spécialisé (ELECTRA) et un LLM généraliste (GPT-4o) permet d'obtenir le
meilleur rapport coût/performance, ouvrant une piste d'amélioration pour des travaux futurs.

\subsubsection{Web scraping et intelligence concurrentielle}

La littérature académique portant spécifiquement sur le web scraping appliqué à l'intelligence
concurrentielle reste comparativement restreinte. Un chapitre publié chez Springer Nature (2025)
\cite{springer2025} propose un état de l'art accompagné d'une étude de cas démontrant que l'intégration du
web scraping dans une stratégie d'entreprise améliore la compréhension du marché ; les auteurs soulignent
également les défis inhérents à cette pratique (structuration des données hétérogènes, maintenance des
extracteurs face aux évolutions des sites) et évoquent la contribution potentielle de l'intelligence
artificielle pour renforcer ces systèmes --- un argument central pour justifier le recours à un extracteur
piloté par IA plutôt qu'à des sélecteurs CSS statiques dans ce projet. Une revue bibliométrique de 2025,
fondée sur l'analyse de 567 publications \cite{bibliometrique2025}, confirme la croissance rapide mais
récente du champ de recherche combinant web scraping et modèles d'IA pour des applications d'entreprise, ce
qui situe ce projet dans un domaine de recherche encore émergent académiquement, malgré une pratique déjà
répandue dans l'industrie.

\subsection{Études webographiques}

En complément des sources académiques, plusieurs ressources techniques non évaluées par les pairs ont été
consultées pour leur pertinence opérationnelle immédiate. Des livres blancs et articles de blog d'éditeurs
spécialisés en extraction de données (Grepsr, Firecrawl) détaillent les bonnes pratiques de collecte pour
l'intelligence concurrentielle : respect des fichiers robots.txt, limitation de la fréquence des requêtes,
exclusion des contenus authentifiés, et constitution d'un historique longitudinal plutôt que de simples
relevés ponctuels --- principes directement appliqués dans la conception du pipeline de collecte de ce projet
(tables d'historique en append-only, seuils de délai entre requêtes).

Sur le plan des outils, les dépôts GitHub des frameworks LangChain et LlamaIndex constituent des références
techniques de facto pour l'implémentation de systèmes RAG, bien que non publiés dans un cadre académique
formel. L'outil ScrapeGraphAI, utilisé pour la collecte de données de ce projet, illustre une tendance
récente au sein de l'écosystème open source consistant à piloter l'extraction web par des modèles de langage
plutôt que par des sélecteurs codés en dur, réduisant la maintenance nécessaire face aux changements de
structure des pages sources. Enfin, une revue bibliométrique et systématique publiée sur F1000Research
(Yesuf et al., 2025) \cite{yesuf2025} portant sur l'adoption de l'IA dans les PME situe le contexte
économique de ce projet : les auteurs identifient l'aide à la décision entrepreneuriale et la perception du
marché comme des axes majeurs de la littérature récente sur la compétitivité des PME assistée par IA, ce qui
correspond précisément au cas d'usage traité ici (une agence de développement web/mobile cherchant à mieux
positionner son offre commerciale).

\subsection{Tableau comparatif des approches}

Le tableau~\ref{tab:comparatif_etat_art} synthétise les approches examinées au regard du besoin du projet
(veille concurrentielle locale, corpus métier, contraintes de stage).

\begin{table}[H]
\centering
\begingroup
\setstretch{1.1}
\small
\begin{tabularx}{\textwidth}{|p{2.6cm}|p{2.8cm}|p{3.2cm}|X|}
\hline
\textbf{Approche} & \textbf{Forces} & \textbf{Limites pour ce projet} & \textbf{Décision} \\
\hline
LLM seul (sans retrieval)
& Réponses fluides en langage naturel
& Hallucinations ; pas d'ancrage sur tarifs/avis réels
& Écarté comme solution principale \\
\hline
RAG générique (Lewis ; Wu ; Yu)
& Ancrage documentaire ; métriques retrieval / fidélité
& Benchmarks souvent génériques (Wikipedia, HotpotQA), peu de corpus tarifaires métier
& Retenu, adapté au corpus concurrentiel \\
\hline
RAG correctif / filtres (Yan et al.)
& Vérifie la pertinence avant génération
& Complexité supplémentaire
& Inspiré : seuil $\tau$, filtre géographique structurel \\
\hline
Sentiment supervisé (BERT / FinBERT)
& Bonnes performances si données annotées
& Coût d'annotation et d'entraînement incompatibles avec le stage
& Écarté en production ; annotation réservée à l'évaluation \\
\hline
Sentiment zero-shot LLM (Simmering ; Beno)
& Pas de dataset labellisé pour entraîner ; extraction structurée possible
& Fiabilité à valider ; pas une probabilité calibrée
& Retenu pour enrichir les fiches \\
\hline
Scraping CSS statique
& Simple, contrôlable
& Cassant sur sites JS dynamiques
& Écarté \\
\hline
Scraping piloté IA (ScrapeGraphAI)
& Résilient aux changements de structure
& Coût API ; variabilité à nettoyer
& Retenu pour la collecte \\
\hline
\end{tabularx}
\endgroup
\caption{Comparaison des approches de la littérature et choix retenus pour le projet}
\label{tab:comparatif_etat_art}
\end{table}

\textbf{Manques identifiés (gaps).}
 (i)~Peu d'évaluations RAG sur des corpus tarifaires / réputationnels
d'agences locales. (ii)~Absence de consensus sur des métriques métier pour l'aide commerciale.
(iii)~Littérature scraping + intelligence concurrentielle encore dominée par la littérature grise.
(iv)~Validation humaine du sentiment zero-shot souvent absente des retours d'expérience industriels
publiés. Ces manques motivent le protocole retenu : corpus réel ICHTUS, séparation développement/test,
diagnostics de retrieval et grille humaine.

\subsection{Analyse critique des approches existantes}

Trois constats critiques se dégagent de cette revue et ont directement orienté les choix techniques du
projet.

Premièrement, concernant le RAG, la littérature d'évaluation \cite{yu2024} souligne l'absence de consensus
sur des bancs d'essai (benchmarks) standardisés, la plupart des travaux étant validés sur des jeux de
données de questions-réponses génériques (HotpotQA, Wikipedia) plutôt que sur des corpus métier réels et
propriétaires. Les enjeux de robustesse et de sécurité restent également des sujets de recherche ouverts,
comme le montrent les travaux sur les vulnérabilités adversariales \cite{xue2024}.

Deuxièmement, concernant l'analyse de sentiment, la littérature converge sur un arbitrage coût/performance
clair : le fine-tuning d'un modèle spécialisé reste supérieur en performance brute mais nécessite un jeu de
données annoté et un coût d'entraînement/inférence significativement plus élevé \cite{simmering2023}, tandis
que l'approche zero-shot par LLM offre une performance déjà satisfaisante pour des tâches de classification
simples sans nécessiter de données annotées, au prix d'une fiabilité moindre sur des cas nuancés (synthèse
SentiEval, citée dans la revue de 2026).

Troisièmement, concernant le web scraping pour l'intelligence concurrentielle, le champ reste dominé par la
littérature grise (blogs d'éditeurs, contenus commerciaux) plutôt que par la recherche académique rigoureuse,
la seule étude de cas académique identifiée \cite{springer2025} étant elle-même récente et signalant
explicitement des défis non résolus (structuration des données, maintenance).

\subsection{Positionnement du projet par rapport à la littérature}

Ce projet se positionne à l'intersection de trois domaines rarement combinés dans un même système au sein de
la littérature consultée : l'extraction automatisée de données web pour l'intelligence concurrentielle,
l'interrogation en langage naturel de données métier via une architecture RAG, et l'analyse de sentiment
zero-shot sur des avis clients non structurés. Ce positionnement répond directement à plusieurs limites
identifiées dans l'analyse critique ci-dessus.

Concernant le RAG, la littérature évalue majoritairement les systèmes sur des corpus génériques ; ce projet
contribue, à l'échelle d'un mémoire de licence, une évaluation empirique sur un corpus métier réel et
propriétaire (données tarifaires et avis clients d'agences de développement web/mobile), avec un protocole
distinguant un jeu de développement (réglage) et un jeu de test indépendant (mesure finale), construit
spécifiquement pour ce domaine (questions pertinentes, hors-sujet, vérité terrain nominative, grille humaine).
La calibration empirique du seuil de confiance, réalisée en observant la distribution réelle des scores de
similarité plutôt qu'en fixant une valeur arbitraire, rejoint la recommandation de la littérature
d'évaluation RAG \cite{yu2024} d'adapter les métriques au contexte réel plutôt que d'appliquer un seuil
universel.

Le filtre géographique structurel ajouté en amont de la génération répond directement à une limite
documentée dans la littérature sur le RAG correctif \cite{yan2024} : la fiabilité du respect d'une consigne
de prompt par un LLM n'est pas garantie, ce qui justifie l'ajout d'une vérification indépendante du modèle
générateur plutôt qu'une confiance aveugle dans l'instruction textuelle --- une démarche méthodologique
directement inspirée de cette famille de travaux.

Concernant l'analyse de sentiment, le choix d'une approche zero-shot par LLM plutôt qu'un modèle fine-tuné
s'appuie explicitement sur l'arbitrage documenté par Simmering et Huoviala (2023) \cite{simmering2023} : en
l'absence d'un jeu de données annoté et dans un contexte de stage à ressources limitées, l'approche zero-shot
constitue un choix raisonné plutôt qu'une solution de facilité, dont les limites (absence de métriques de
précision formelles, performance potentiellement inférieure sur des cas nuancés) sont explicitement
documentées dans ce mémoire plutôt que dissimulées.

Enfin, ce projet illustre concrètement la piste évoquée par le chapitre Springer (2025) \cite{springer2025}
consistant à renforcer les systèmes de web scraping par l'intelligence artificielle : l'utilisation d'un
extracteur piloté par LLM (ScrapeGraphAI) plutôt que par des sélecteurs statiques constitue une application
directe de cette recommandation, avec un bénéfice observé en pratique sur la résilience du pipeline face aux
évolutions de structure des pages sources.
